% Options for packages loaded elsewhere
\PassOptionsToPackage{unicode}{hyperref}
\PassOptionsToPackage{hyphens}{url}
\documentclass[
  11pt,a4paper]{article}
\usepackage{xcolor}
\usepackage[margin=1in]{geometry}
\usepackage{amsmath,amssymb}
\setcounter{secnumdepth}{-\maxdimen} % remove section numbering
\usepackage{iftex}
\ifPDFTeX
  \usepackage[T1]{fontenc}
  \usepackage[utf8]{inputenc}
  \usepackage{textcomp} % provide euro and other symbols
\else % if luatex or xetex
  \usepackage{unicode-math} % this also loads fontspec
  \defaultfontfeatures{Scale=MatchLowercase}
  \defaultfontfeatures[\rmfamily]{Ligatures=TeX,Scale=1}
\fi
\usepackage{lmodern}
\ifPDFTeX\else
  % xetex/luatex font selection
  \setmainfont[]{Latin Modern Roman}
\fi
% Use upquote if available, for straight quotes in verbatim environments
\IfFileExists{upquote.sty}{\usepackage{upquote}}{}
\IfFileExists{microtype.sty}{% use microtype if available
  \usepackage[]{microtype}
  \UseMicrotypeSet[protrusion]{basicmath} % disable protrusion for tt fonts
}{}
\makeatletter
\@ifundefined{KOMAClassName}{% if non-KOMA class
  \IfFileExists{parskip.sty}{%
    \usepackage{parskip}
  }{% else
    \setlength{\parindent}{0pt}
    \setlength{\parskip}{6pt plus 2pt minus 1pt}}
}{% if KOMA class
  \KOMAoptions{parskip=half}}
\makeatother
\usepackage{longtable,booktabs,array}
\newcounter{none} % for unnumbered tables
\usepackage{calc} % for calculating minipage widths
% Correct order of tables after \paragraph or \subparagraph
\usepackage{etoolbox}
\makeatletter
\patchcmd\longtable{\par}{\if@noskipsec\mbox{}\fi\par}{}{}
\makeatother
% Allow footnotes in longtable head/foot
\IfFileExists{footnotehyper.sty}{\usepackage{footnotehyper}}{\usepackage{footnote}}
\makesavenoteenv{longtable}
\setlength{\emergencystretch}{3em} % prevent overfull lines
\providecommand{\tightlist}{%
  \setlength{\itemsep}{0pt}\setlength{\parskip}{0pt}}
\usepackage{mathtools}
\usepackage{enumitem}
\usepackage{microtype}
\usepackage{booktabs}
\usepackage{array}
\usepackage{bookmark}
\IfFileExists{xurl.sty}{\usepackage{xurl}}{} % add URL line breaks if available
\urlstyle{same}
\hypersetup{
  pdftitle={Foundations of Substrate Gravity From Newtonian Limit to Curvature Emergence in the Event Density Framework},
  pdfauthor={Allen Proxmire},
  hidelinks,
  pdfcreator={LaTeX via pandoc}}

\title{Foundations of Substrate Gravity\\
\large From Newtonian Limit to Curvature Emergence in the Event Density
Framework}
\author{Allen Proxmire}
\date{April 2026}

\begin{document}
\maketitle

\subsection{Abstract}\label{abstract}

The Event Density (ED) program presents a substrate ontology in which
classical gravitation emerges as a structural consequence of substrate
participation-channel dynamics rather than as a fundamental field
theory. This paper consolidates the closed structural-foundation work of
the program's substrate-gravity sector into a single
architectural-and-substrate account spanning flat-background regime
through curvature emergence. \textbf{Theorem T19} derives Newton's
gravitational constant \(G = c^3\ell_P^2/\hbar\) from substrate
participation imbalance + cumulative-strain mechanism + holographic
participation-count bound. \textbf{Theorem T20} derives the transition
acceleration \(a_0 = c\,H_0/(2\pi)\) from cosmological-horizon
participation density + retardation + substrate spin-statistics +
holographic bound; the result is \(\ell_P\)-invariant under continuum
limit, structurally distinct from substrate-cutoff-anchored mass-gap
mechanisms in other ED sectors. The \textbf{ED Combination Rule}
\(a = \sqrt{a_N\,a_0}\) supplies the substrate-level geometric-mean
composition between Newton-class and transition-class accelerations.
\textbf{Theorem T21} derives the baryonic Tully-Fisher relation
\(v^4 = G\,M\,a_0\) as a substrate consequence of T19 + T20 + ECR. Under
DCGT (Diffusion Coarse-Graining Theorem) the substrate-gravity content
consolidates into the modified Poisson equation
\(\nabla\cdot[\mu(|\nabla\Phi|/a_0)\nabla\Phi] = 4\pi G\rho_m\) with
FORCED asymptotic constraints; BTFR slope-4 is robust under all
admissible kernel-profile / interpolation-function /
substrate-discreteness variations.

The curvature-emergence extension (Arc ED-10) identifies the substrate
cumulative-strain four-index object as the load-bearing curvature degree
of freedom and produces an acoustic-metric-class scalar-tensor
covariantization
\(\nabla_\mu[\mu(\sqrt{g^{\alpha\beta}\nabla_\alpha\Phi\nabla_\beta\Phi}/a_0)\nabla^\mu\Phi] = 4\pi G\,T\)
as the substrate-FORCED covariant generalization. The covariant equation
reduces to the flat-background modified Poisson equation in the
weak-field limit; it preserves OS positivity + ghost-freedom + gradient
stability under structural conditions; deep-MOND superluminality is
structurally inherited from the substrate-derivation chain as the only
structural cost. The verdict is \textbf{conditional-positive at
structural-suggestive level} parallel in form and honesty to the
program's NS-Smoothness Intermediate Path C and Yang-Mills
Clay-relevance verdicts. \textbf{No claim of general relativity
emergence is made}; the Einstein equation
\(G_{\mu\nu} = 8\pi G\,T_{\mu\nu}\) remains SPECULATIVE per ED-Phys-10
acoustic-metric-only baseline. What ED \emph{does} deliver is a
substrate-grounded architectural account of Newtonian gravity +
MOND-class behavior + BTFR slope-4 + acoustic-metric curvature emergence
with explicit form-FORCED / value-INHERITED demarcation. The paper is
the canonical substrate-gravity-and-curvature-emergence reference for
the ED program.

\begin{center}\rule{0.5\linewidth}{0.5pt}\end{center}

\subsection{1. Introduction}\label{introduction}

\subsubsection{1.1 Motivation}\label{motivation}

Classical gravitation is presented in standard physics as either
Newton's law of universal gravitation (a phenomenological inverse-square
attractive force) at solar-system scales, or as Einstein's general
relativity (a fundamental dynamical field theory of the spacetime
metric) at relativistic scales. The MOND-class modification ---
Milgrom's hypothesis \(\mu(|\mathbf{a}|/a_0)\mathbf{a} = \mathbf{a}_N\)
at the level of acceleration scaling --- was introduced
phenomenologically to explain galactic rotation curves without invoking
dark matter, and has been extensively studied as both a phenomenological
framework and via various covariantization attempts (RAQUAL,
scalar-tensor MOND, full TeVeS, etc.) since Bekenstein-Milgrom 1984. The
empirical regularity of the baryonic Tully-Fisher relation
\(v^4 \propto M\) across galaxy samples is one of the strongest
empirical anchors of the MOND-class framework; standard physics has no
first-principles derivation of the slope-4 power.

The Event Density (ED) program supplies a substrate ontology under which
Newton's gravitational coupling, the MOND transition acceleration
\(a_0\), the geometric-mean composition rule between Newton-class and
transition-class regimes, the BTFR slope-4 relation, and the
modified-Poisson-class field equation all emerge as substrate
consequences of a single set of substrate primitives. The goal is not to
replace general relativity at the constructive-rigorous level, but to
identify the substrate-derivation chain that produces the leading-order
classical-gravitation content + its empirically-anchored MOND-class
extensions.

\subsubsection{1.2 Relation to ED Program}\label{relation-to-ed-program}

The substrate-gravity content sits within the broader Event Density
structural-foundation work that includes Phase-1 QM emergence (sixteen
forced theorems for the four QM postulates), the closed
structural-foundation theorems T17 (gauge-fields-as-rule-type), T18 (V1
kernel retardation), Theorem N1 (V1 finite-width vacuum kernel), Theorem
GR1 (V1 with Synge world function on curved spacetime), and the
substrate-to-continuum bridge theorem DCGT (Arc D, 2026-04-30). The
substrate-gravity work draws on T18 + Theorem N1 + GR1 + DCGT directly;
the form-FORCED / value-INHERITED methodology is consistent across every
ED program sector.

The substrate-gravity content has been developed across three closed
arcs: (a) the original substrate-gravity arc (T19 / T20 / ECR / T21,
2026-04-27/28), (b) the Substrate Gravity Extension Arc (Arc SG,
2026-04-30) consolidating the closed substrate-gravity work under DCGT
methodology + producing the modified Poisson equation, (c) the Curvature
Emergence Arc (Arc ED-10, 2026-04-30) extending to acoustic-metric-class
curvature emergence + producing the scalar-tensor covariantization. The
present paper integrates all three arcs into a single canonical
reference.

\subsubsection{1.3 Scope}\label{scope}

The paper covers the flat-background regime + curvature-emergent regime
under the acoustic-metric-class baseline of ED-Phys-10. Scope items:

\begin{itemize}
\tightlist
\item
  Substrate origin of Newton's \(G\) and the transition acceleration
  \(a_0\).
\item
  Substrate origin of the geometric-mean composition rule.
\item
  BTFR slope-4 derivation + robustness audit.
\item
  Modified Poisson equation as flat-background substrate-gravity field
  equation.
\item
  Curvature-emergent scalar-tensor covariantization.
\item
  OS positivity + ghost-freedom + gradient-stability audit at
  structural-suggestive level.
\end{itemize}

\subsubsection{1.4 Non-goals}\label{non-goals}

Explicitly out of scope:

\begin{itemize}
\tightlist
\item
  \textbf{No claim of deriving general relativity.} The Einstein
  equation \(G_{\mu\nu} = 8\pi G T_{\mu\nu}\) remains SPECULATIVE per
  ED-Phys-10 acoustic-metric-only baseline.
\item
  \textbf{No new primitives.} The thirteen-primitive substrate framework
  is preserved throughout.
\item
  \textbf{No cosmology or dark-energy modeling beyond what \(a_0\)
  supplies.} Specific cosmological-scale physics (vacuum energy,
  dark-energy equation of state, structure formation) out of scope.
\item
  \textbf{No metric dynamics beyond what substrate structure forces.}
  Gravitational waves, black-hole solutions, cosmological metric
  evolution are addressed only insofar as substrate-derivation FORCES
  specific structural content.
\item
  \textbf{No claim of solving any Clay Millennium Problem.} Curvature
  emergence is not itself a Clay problem; the structural-positive
  content is the substrate-grounding of MOND-class behavior.
\end{itemize}

\begin{center}\rule{0.5\linewidth}{0.5pt}\end{center}

\subsection{2. Substrate Primitives and Participation
Structure}\label{substrate-primitives-and-participation-structure}

The ED substrate consists of a network of \emph{chains} --- fundamental
discrete entities that carry participation channels between
substrate-level events. Chains support participation channels via
finite-bandwidth commitments at substrate-level events; the substrate
dynamics is governed by thirteen primitives (P-01 through P-13) that fix
the structural content without committing to numerical values at
substrate level.

\textbf{V1 kernel} --- the vacuum-response participation kernel ---
supplies the substrate-level finite-width spatial smoothing structure.
By Theorem N1 (Arc N closure), the V1 kernel has finite spatial width
\(\sim\ell_P\), smooth profile, isotropic by substrate-level rotational
invariance, and positive Fourier transform. By Theorem 18 (Arc B
closure), the V1 kernel has forward-cone-only support, anchoring the
kernel-level arrow of time at substrate level.

\textbf{Participation density} --- the substrate-level
participation-channel density at a given spacetime point --- is the
scalar-field-class quantity that maps under DCGT (Arc D closure) to the
macroscopic mass density at fluid-mechanical / gravitational-substrate
scale. Variations in participation density source the gravitational
potential gradient via the acoustic-metric reading of Arc ED-10.

\textbf{Directional bias} --- the substrate-level orientation-bearing
participation structure --- supplies the directional-field-class
quantities that map to vector fields (velocity, gauge potentials,
magnetic field) at continuum scale. Per ED-I-06 (Fields and Forces, Feb
2026), the three field classes (directional, scalar, curvature-like) are
exhaustive at substrate level.

\textbf{Cumulative strain} --- the substrate-level multi-point
participation-channel correlation that captures how parallel transport
across multi-point substrate loops deviates from triviality --- is the
load-bearing substrate quantity for both T19's substrate-gravity
mechanism and Arc ED-10's curvature-emergent degrees of freedom. The
four-point cumulative-strain object supplies the substrate origin of
Riemann curvature at continuum scale.

\textbf{DCGT (Diffusion Coarse-Graining Theorem)} --- the
substrate-to-continuum bridge theorem closed in Arc D --- establishes
the methodology for deriving continuum-scale equations from substrate
flux statistics under hydrodynamic-window scale separation
\(\ell_P \ll R_\mathrm{cg} \ll L_\mathrm{flow}\). Multi-scale expansion
of substrate fluxes produces continuum-scale operators with form-FORCED
structure + value-INHERITED coefficients. DCGT covers scalar diffusion
(mobility-channel), directional-field viscosity, V1-second-moment
hyperviscous corrections, V5 cross-chain temporal-memory contributions,
T17-minimal-coupling matter-source content, and (per Arc ED-10 §3)
generalized to curvature-like-field-class content via Laplace-Beltrami
operator + V1-second-moment curvature-coupling terms.

These five structural ingredients (V1 kernel, participation density,
directional bias, cumulative strain, DCGT) form the substrate basis of
the substrate-gravity work that follows.

\begin{center}\rule{0.5\linewidth}{0.5pt}\end{center}

\subsection{3. Theorem T19 --- Newton Constant from
Substrate}\label{theorem-t19-newton-constant-from-substrate}

\subsubsection{Theorem statement}\label{theorem-statement}

\begin{quote}
\textbf{Theorem T19 (Newton's gravitational constant from substrate,
FORCED-unconditional).} \emph{Under ED substrate primitives +
cumulative-strain mechanism + holographic participation-count bound +
dimensional-analysis closure, Newton's gravitational constant takes the
form}

\[G \;=\; \frac{c^3\,\ell_P^2}{\hbar}.\]

\emph{The form is FORCED at substrate level by the
dimensional-analysis-tight combination of substrate length scale
\(\ell_P\), substrate-level light-speed scale \(c\), and substrate-level
participation-quantum scale \(\hbar\). The numerical value of \(G\) is
INHERITED at value layer through \(\ell_P\), \(c\), \(\hbar\) as
substrate inputs.}
\end{quote}

\subsubsection{Derivation sketch}\label{derivation-sketch}

The substrate-derivation chain runs through three structural anchors:

\textbf{(a) Substrate participation imbalance.} Mass content \(M\) at a
substrate position sources a participation-density bias that propagates
outward through the chain network with \(1/r^2\) geometric dilution (3D
spherical isotropy). The participation-flux density at radius \(r\) is
set by the total mass-content modulated by substrate
participation-conservation.

\textbf{(b) Cumulative-strain mechanism.} The substrate-level cumulative
strain across a substrate region of size \(r\) accumulates from the
participation imbalance integrated over the region; the strain-content
scales with the holographic participation-count bound at the region's
boundary.

\textbf{(c) Holographic participation-count bound + dimensional
analysis.} The participation-count at scale \(r\) is bounded as
\(N \sim r^2/\ell_P^2\) (holographic-class). Combining this with the
substrate-level proper-time scale (\(\sim \ell_P/c\)) and the
substrate-level participation-quantum (\(\hbar\)) gives the
dimensionally-tight combination producing \(G\).

The detailed derivation closure is in \texttt{arcs/arc-SG/}; the result
is FORCED at substrate level with value INHERITED via \(\ell_P\) + \(c\)
+ \(\hbar\).

\subsubsection{Structural meaning}\label{structural-meaning}

Newton's gravitational coupling is not a fundamental constant in the ED
framework; it is a substrate-derived combination of more fundamental
substrate scales. The \(\ell_P^2\) scaling identifies \(G\) as a
substrate-cutoff-anchored quantity: variations of the substrate length
scale \(\ell_P\) produce proportional variations in \(G\) at value
layer. The substrate-derivation supplies a structural reason why \(G\)
has its specific dimensional form --- it is the unique combination of
substrate scales producing a quantity with \(G\)'s dimensions.

This positions Newton's gravity as a substrate-structural feature rather
than a phenomenologically-fitted coupling. Standard physics treats \(G\)
as an empirical constant; the ED substrate-derivation supplies its
substrate origin while preserving the empirical-input nature of the
numerical value.

\begin{center}\rule{0.5\linewidth}{0.5pt}\end{center}

\subsection{\texorpdfstring{4. Theorem T20 --- Transition Acceleration
\(a_0\)}{4. Theorem T20 --- Transition Acceleration a\_0}}\label{theorem-t20-transition-acceleration-a_0}

\subsubsection{Theorem statement}\label{theorem-statement-1}

\begin{quote}
\textbf{Theorem T20 (Transition acceleration from substrate,
FORCED-unconditional).} \emph{Under ED substrate primitives +
cosmological-horizon participation density + T18 retardation + substrate
spin-statistics + holographic bound at horizon scale, the transition
acceleration takes the form}

\[a_0 \;=\; \frac{c\,H_0}{2\pi}.\]

\emph{The form is FORCED at substrate level + the \(2\pi\) prefactor is
structurally derived from substrate spin-statistics + holographic-bound
normalizations. The numerical value of \(a_0\) depends on the
cosmological Hubble parameter \(H_0\), INHERITED at value layer.}
\end{quote}

\subsubsection{\texorpdfstring{\(\ell_P\)-invariance under continuum
limit}{\textbackslash ell\_P-invariance under continuum limit}}\label{ell_p-invariance-under-continuum-limit}

Per Arc\_SG\_3 (Substrate Gravity Extension), \(a_0\) is
\textbf{invariant under \(\ell_P \to 0\)}. The substrate-derivation
routes through three anchors all preserved under substrate-cutoff scale
variation:

\textbf{(A1) Cosmological-horizon participation density.}
\(R_H = c/H_0\) is independent of substrate-cutoff scale. The
horizon-scale participation density depends on cosmological parameters
(\(H_0\), vacuum-energy content) which are independent of \(\ell_P\).

\textbf{(A2) Retardation structure (T18).} Forward-cone-only support of
V1 kernel is FORCED at primitive level (Arc B closure); preservation
under \(\ell_P \to 0\) structural.

\textbf{(A3) \(2\pi\) prefactor.} Substrate-spin-statistics +
holographic-bound normalizations cancel the \(\ell_P\)-dependence
(\(N \sim R_H^2/\ell_P^2\) ×
substrate-spin-statistics-normalization-with-\(\ell_P^{-2}\)-suppression).
The prefactor is dimensionless + \(\ell_P\)-independent.

\subsubsection{Horizon-scale participation
density}\label{horizon-scale-participation-density}

The cosmological event horizon at radius \(R_H = c/H_0\) supplies a
substrate-level participation-density boundary. Beyond \(R_H\), no
substrate participation channels can communicate with the local
participation network within the cosmological time available; the
horizon truncates substrate participation-flux integration at scale
\(R_H\). The horizon-truncation at scale \(R_H\) produces the
transition-regime acceleration scale \(a_0 \sim cH_0\) via dimensional
analysis + substrate-derivation chain.

\subsubsection{Contrast with YM mass-gap
scaling}\label{contrast-with-ym-mass-gap-scaling}

The substrate origin of \(a_0\) is fundamentally different from the
substrate origin of mass gaps in other ED sectors. Specifically,
contrast with the Yang-Mills mass-gap mechanism (Arc YM-3):

{\def\LTcaptype{none} % do not increment counter
\begin{longtable}[]{@{}
  >{\raggedright\arraybackslash}p{(\linewidth - 4\tabcolsep) * \real{0.3333}}
  >{\raggedright\arraybackslash}p{(\linewidth - 4\tabcolsep) * \real{0.3333}}
  >{\raggedright\arraybackslash}p{(\linewidth - 4\tabcolsep) * \real{0.3333}}@{}}
\toprule\noalign{}
\begin{minipage}[b]{\linewidth}\raggedright
Substrate quantity
\end{minipage} & \begin{minipage}[b]{\linewidth}\raggedright
Substrate scale
\end{minipage} & \begin{minipage}[b]{\linewidth}\raggedright
Continuum-limit survival
\end{minipage} \\
\midrule\noalign{}
\endhead
\bottomrule\noalign{}
\endlastfoot
YM mass gap \(m_\mathrm{eff}^2\) & \(\ell_P^{-2}\) (substrate cutoff,
UV) & \textbf{Conditional} on kernel-profile rescaling \\
\(a_0\) transition acceleration & \(cH_0\) (cosmological horizon, IR) &
\textbf{Automatic}; no rescaling required \\
\end{longtable}
}

The YM mass gap is a \emph{substrate-cutoff effect} (UV); \(a_0\) is a
\emph{cosmological-horizon effect} (IR). YM mass gap requires
kernel-profile rescaling for continuum-limit survival; \(a_0\) inherits
no \(\ell_P\)-dependence and survives automatically. This structural
distinction is one of the key findings of the substrate-gravity arc.

\begin{center}\rule{0.5\linewidth}{0.5pt}\end{center}

\subsection{5. ECR --- ED Combination
Rule}\label{ecr-ed-combination-rule}

\subsubsection{Statement}\label{statement}

The ED Combination Rule (ECR), closed in
\texttt{arcs/arc-SG/ED\_combination\_rule.md}, states that the magnitude
of the substrate-gravitational acceleration is the geometric mean of the
Newton-class and transition-class accelerations in the deep-MOND regime:

\[a \;=\; \sqrt{a_N\,a_0}, \qquad a_N \;=\; |\nabla\Phi_N|,\]

where \(\Phi_N\) is the Newtonian potential satisfying the standard
Poisson equation. Squaring,

\[a^2 \;=\; a_N\,a_0,\]

at the magnitude level of the substrate-gravitational acceleration in
the deep-MOND regime.

\subsubsection{Substrate origin}\label{substrate-origin}

ECR is FORCED at substrate level by the participation-curvature
combination at the cross-term scale between Newton-class
(substrate-cutoff-anchored) and transition-class
(cosmological-horizon-anchored) substrate sectors. The
substrate-derivation routes through the cross-term
\(\sqrt{GMa_0}\cdot\log(R/R_0)\) in the deep-MOND regime; the
geometric-mean composition reflects the substrate-level
participation-curvature interpolation between the two regimes.

\subsubsection{Role in deep-MOND
behavior}\label{role-in-deep-mond-behavior}

ECR's geometric-mean composition is the substrate origin of MOND-class
behavior. Specifically:

\begin{itemize}
\tightlist
\item
  In the high-acceleration regime (\(a_N \gg a_0\)), the
  substrate-gravity acceleration approaches \(a_N\) --- Newton's law is
  recovered.
\item
  In the deep-MOND regime (\(a_N \ll a_0\)), the substrate-gravity
  acceleration approaches \(\sqrt{a_N\,a_0}\) --- the substrate-level
  MOND-class behavior.
\item
  The crossover occurs at acceleration scale \(a_0\).
\end{itemize}

ECR supplies the structural framework within which the modified Poisson
equation (Section 7) emerges as the field-equation-level expression of
the substrate composition.

\subsubsection{Naming convention}\label{naming-convention}

Per program-level memory: the substrate rule is called ``the ED
Combination Rule'' (ECR) --- no number assigned. ``P11'' is reserved for
Primitive 11 (commitment-irreversibility); ``ED-11'' should not be used
in naming.

\begin{center}\rule{0.5\linewidth}{0.5pt}\end{center}

\subsection{6. Theorem T21 --- BTFR
Slope-4}\label{theorem-t21-btfr-slope-4}

\subsubsection{Derivation}\label{derivation}

In the deep-MOND regime, ECR supplies \(a^2 = a_N\,a_0\) with
\(a_N = GM/r^2\) for spherical mass distribution. Substituting:

\[a^2 \;=\; \frac{G\,M\,a_0}{r^2} \quad\Longrightarrow\quad |a(r)| \;=\; \frac{\sqrt{G\,M\,a_0}}{r}.\]

For a circular orbit at radius \(r\), centripetal balance gives
\(v^2/r = |a(r)|\), hence \(v^2 = \sqrt{G\,M\,a_0}\). Squaring,

\[\boxed{\;v^4 \;=\; G\,M\,a_0.\;}\]

This is the \textbf{baryonic Tully-Fisher relation} at slope-4 ---
Theorem T21.

\subsubsection{Why slope-4 is FORCED}\label{why-slope-4-is-forced}

Slope-4 is FORCED by:

\begin{itemize}
\tightlist
\item
  \textbf{ECR's geometric-mean composition} at substrate level (FORCED
  via closed-arc derivation).
\item
  \textbf{Spherical symmetry} of test-mass distribution (assumed for
  galactic-rotation-curve idealization at large radii).
\item
  \textbf{Circular-orbit centripetal balance} (kinematic).
\item
  \textbf{Dimensional analysis.} The only dimensional combination of
  \(G\), \(M\), and \(a_0\) producing {[}velocity{]}⁴ is \(GMa_0\).
\end{itemize}

Each ingredient is structurally backed; their combination yields slope-4
as a substrate consequence.

\subsubsection{Robustness}\label{robustness}

Per Arc\_SG\_5 (BTFR robustness audit), slope-4 is \textbf{structurally
robust} under all admissible variations:

\begin{itemize}
\tightlist
\item
  V1 kernel-profile variations (within Theorem N1 admissibility class).
\item
  Interpolation-function \(\mu(x)\) variations (within
  asymptotic-constraint class --- Section 7).
\item
  Substrate-discreteness variations (within DCGT hydrodynamic-window
  conditions).
\item
  Coarse-graining-window variations (within hydrodynamic window).
\end{itemize}

ED substrate-gravity predicts \textbf{zero intrinsic scatter} in the
deep-MOND asymptotic regime. Observed BTFR scatter (sub-0.1 dex in SPARC
samples; slope \(\sim 3.85 \pm 0.15\)) is dominated by baryonic-mass
measurement uncertainty + transition-regime \(\mu(x)\) + baryonic
geometry. This is consistent with empirical galaxy-rotation-curve data
and supplies an empirical anchor for the substrate-derivation chain.

\begin{center}\rule{0.5\linewidth}{0.5pt}\end{center}

\subsection{7. Arc SG Synthesis --- Flat-Background Substrate
Gravity}\label{arc-sg-synthesis-flat-background-substrate-gravity}

\subsubsection{Consolidated field
equation}\label{consolidated-field-equation}

Under DCGT methodology, the closed substrate-gravity work consolidates
into the modified Poisson equation in flat background:

\[\boxed{\;\nabla\cdot\!\left[\mu\!\left(\frac{|\nabla\Phi|}{a_0}\right)\nabla\Phi\right] \;=\; 4\pi G\,\rho_m,\;}\]

with substrate-derived gravitational coupling \(G = c^3\ell_P^2/\hbar\)
(T19), substrate-derived transition acceleration \(a_0 = c\,H_0/(2\pi)\)
(T20, \(\ell_P\)-invariant), and dimensionless interpolation function
\(\mu(x)\) satisfying FORCED asymptotic constraints.

\subsubsection{Newtonian limit}\label{newtonian-limit}

In the high-acceleration limit (\(|\nabla\Phi| \gg a_0\), equivalently
\(x \gg 1\)), \(\mu(x) \to 1\) and the field equation reduces to the
standard Newtonian Poisson equation:

\[\nabla^2\Phi \;=\; 4\pi G\,\rho_m.\]

Newton's law of gravitation \(\Phi(r) = -GM/r\) + force law
\(\mathbf{F} = -GMm/r^2\,\hat r\) are recovered as standard 3D
Green's-function solutions.

\subsubsection{Deep-MOND limit}\label{deep-mond-limit}

In the low-acceleration regime (\(|\nabla\Phi| \ll a_0\), equivalently
\(x \ll 1\)), \(\mu(x) \to x\) and the field equation reduces to:

\[\nabla\cdot\!\left[\frac{|\nabla\Phi|}{a_0}\,\nabla\Phi\right] \;=\; 4\pi G\,\rho_m.\]

For spherical symmetry, the first integral outside the source gives
\(|\nabla\Phi|^2 = GMa_0/r^2\) (Section 6 derivation); circular-orbit
balance produces BTFR slope-4 \(v^4 = GMa_0\).

\subsubsection{BTFR}\label{btfr}

BTFR slope-4 follows from the deep-MOND asymptotic + spherical symmetry
+ circular-orbit balance + dimensional analysis. T21's empirical-anchor
result is recovered as a substrate consequence.

\subsubsection{FORCED vs INHERITED
structure}\label{forced-vs-inherited-structure}

\textbf{FORCED at substrate level:} - Modified Poisson equation form
(DCGT scalar-diffusion + ECR substrate composition). - Asymptotic
constraint \(\mu(x) \to 1\) at \(x \to \infty\) (Newtonian-regime
substrate consistency). - Asymptotic constraint \(\mu(x) \to x\) at
\(x \to 0\) (ECR substrate composition). - BTFR slope-4 (deep-MOND
asymptotic + dimensional analysis). - Reduction to Newtonian Poisson at
high acceleration.

\textbf{INHERITED at value layer:} - Detailed shape of \(\mu(x)\) in
transition regime (multiple structurally-equivalent forms: ``simple''
\(x/\sqrt{1+x^2}\), ``standard'' \(x/(1+x)\), exponential, rational
families). - Numerical values of \(G\) and \(a_0\) (via \(\ell_P\) and
\(H_0\)). - Cosmological boundary conditions; specific mass
distributions.

The flat-background substrate-gravity sector is structurally complete at
the field-equation level under DCGT methodology.

\begin{center}\rule{0.5\linewidth}{0.5pt}\end{center}

\subsection{8. Curvature Emergence (Arc
ED-10)}\label{curvature-emergence-arc-ed-10}

\subsubsection{Curvature-like degrees of
freedom}\label{curvature-like-degrees-of-freedom}

Per Arc\_ED10\_2, the substrate quantities with curvature-like index
structure are:

\begin{itemize}
\tightlist
\item
  \textbf{Cumulative-strain four-index object} --- the load-bearing
  substrate degree of freedom for curvature. The genuine independent
  substrate quantity capturing loop-holonomy-class
  participation-curvature content.
\item
  \textbf{Participation-curvature tensor from GR1} --- diagnostic /
  cross-checking; derived from underlying metric structure.
\item
  \textbf{Directional-bias-gradient commutators} --- secondary; right
  index structure but not curvature-vanishing in flat space.
\item
  \textbf{Kernel-shape second-derivative structure} --- diagnostic;
  responds to curvature, doesn't produce it.
\end{itemize}

The cumulative-strain four-index object is uniquely the substrate origin
of Riemann curvature at continuum scale.

\subsubsection{Acoustic metric}\label{acoustic-metric}

The substrate participation-channel network supplies a metric-like
quantity at continuum scale via the participation-density-modulated
effective metric:

\[g_{\mu\nu}^\mathrm{acoustic}(x) \;\propto\; \bigl(\text{participation-density-modulated effective metric}\bigr).\]

The acoustic-metric reading is the standard ED interpretation per
ED-Phys-10 acoustic-metric-only baseline: the metric is \emph{kinematic}
(a continuum-level summary of substrate participation-density variation)
rather than \emph{dynamical} (an independent field with its own equation
of motion). The Christoffel-class connection follows kinematically; the
Riemann-like curvature object follows kinematically from the metric.

\subsubsection{Weak-field reduction to
SG-4}\label{weak-field-reduction-to-sg-4}

Per Arc\_ED10\_3, the curvature-emergent framework reduces to the
flat-background modified Poisson equation in the
\(g_{\mu\nu}\to\eta_{\mu\nu}\) limit under three load-bearing structural
assumptions:

\begin{itemize}
\tightlist
\item
  \textbf{(A1) Acoustic-metric class} (kinematic metric, not fundamental
  dynamical).
\item
  \textbf{(A2) Levi-Civita connection} (metric-derived kinematically).
\item
  \textbf{(A3) Subleading curvature-coupling coefficients}
  (\(\alpha_R\), \(\beta_{R_{\mu\nu}}\) from curved-DCGT operator
  \(\ell_P^2\)-suppressed at leading weak-field order).
\end{itemize}

Standard GR weak-field-limit identities are preserved:
\(h_{00} \approx 2\Phi/c^2\) FORCED + \(R_{00} \propto \nabla^2\Phi\)
FORCED + curved-DCGT operator reduces to flat Laplacian at leading order
FORCED. All six SG-6 weak-field-limit prerequisites verified.

\subsubsection{Scalar-tensor
covariantization}\label{scalar-tensor-covariantization}

Per Arc\_ED10\_4, three structural families audited against five
substrate-compatibility constraints:

{\def\LTcaptype{none} % do not increment counter
\begin{longtable}[]{@{}
  >{\raggedright\arraybackslash}p{(\linewidth - 2\tabcolsep) * \real{0.5000}}
  >{\raggedright\arraybackslash}p{(\linewidth - 2\tabcolsep) * \real{0.5000}}@{}}
\toprule\noalign{}
\begin{minipage}[b]{\linewidth}\raggedright
Family
\end{minipage} & \begin{minipage}[b]{\linewidth}\raggedright
Status
\end{minipage} \\
\midrule\noalign{}
\endhead
\bottomrule\noalign{}
\endlastfoot
TeVeS-class scalar-tensor (in scalar-tensor reduced form) &
\textbf{Compatible} --- passes all five constraints \\
\(f(R)\)-class & \textbf{Inadmissible} --- violates ED-Phys-10
(introduces scalaron + dynamical metric) \\
Bimetric-class & \textbf{Equivalent to TeVeS-class} under
acoustic-metric identification \\
\end{longtable}
}

The substrate FORCES the scalar-tensor acoustic-metric covariantization
as the unique compatible candidate-form. The result is structurally
equivalent to RAQUAL / scalar-tensor MOND covariantizations in the
standard literature; the ED-program-specific contribution is the
substrate-derivation chain identifying this form as substrate-FORCED.

\subsubsection{OS positivity + ghost-free +
gradient-stable}\label{os-positivity-ghost-free-gradient-stable}

Per Arc\_ED10\_5: - \textbf{OS positivity preserved} under (C1)
monotone-increasing \(\mu(x)\) + (C2) kinetic-matrix positivity
\(\mu(x) + x\mu'(x) > 0\) + (C3) V1 positive Fourier transform. -
\textbf{Ghost-freedom holds} under (C2). Kinetic matrix
positive-definite in transverse + longitudinal directions. -
\textbf{Gradient stability holds} under (C2). No spatial-mode
instabilities at linearized scalar-sector level.

Conditions (C1)--(C3) are structurally backed by closed-arc work; (C1)
and (C2) FORCED in asymptotic limits + standard \(\mu(x)\) forms; (C3)
FORCED unconditionally by Theorem 18 + Theorem N1.

\subsubsection{Deep-MOND superluminality
(structural-conditional)}\label{deep-mond-superluminality-structural-conditional}

Per Arc\_ED10\_5 Step 3: in the deep-MOND regime (\(\mu \to x\),
\(\mu' \to 1\)), the longitudinal scalar-mode propagates superluminally
relative to the acoustic light cone (longitudinal/transverse eigenvalue
ratio = 2; longitudinal cone wider by factor \(\sqrt{2}\)). The
superluminality is \textbf{structurally FORCED by the
substrate-derivation chain} that produces SG-4 + ECR + BTFR slope-4.
Avoidance routes structurally inadmissible:

\begin{itemize}
\tightlist
\item
  Adding an independent vector field (Bekenstein full TeVeS) violates
  the no-new-primitives constraint.
\item
  Modifying \(\mu(x)\) to \(\mu' \to 0\) in deep-MOND violates the BTFR
  slope-4 FORCED-asymptotic.
\end{itemize}

Superluminality is the structural cost of obtaining the
substrate-derived MOND-covariantization. Causality is preserved via
effective \(g_\mathrm{eff}\)-cone consistency in physically reasonable
backgrounds (standard scalar-tensor MOND-covariantization-literature
view).

\begin{center}\rule{0.5\linewidth}{0.5pt}\end{center}

\subsection{9. Covariant Field Equation}\label{covariant-field-equation}

The substrate-FORCED covariant generalization of the modified Poisson
equation (Section 7) is:

\[\boxed{\;\nabla_\mu\!\left[\mu\!\left(\frac{\sqrt{g^{\alpha\beta}\nabla_\alpha\Phi\,\nabla_\beta\Phi}}{a_0}\right)\nabla^\mu\Phi\right] \;=\; 4\pi G\,T,\;}\]

with:

\begin{itemize}
\tightlist
\item
  \(g_{\mu\nu}\) the acoustic metric (Section 8) --- kinematic-summary
  of substrate participation-density variation per ED-Phys-10 baseline.
\item
  \(\nabla_\mu\) the metric-covariant derivative associated with the
  Levi-Civita connection of \(g_{\mu\nu}\).
\item
  \(\mu(x)\) the dimensionless interpolation function with FORCED
  asymptotic constraints from Section 7 (\(\mu(x) \to 1\) at
  \(x \to \infty\), \(\mu(x) \to x\) at \(x \to 0\)).
\item
  \(T\) the trace of the matter stress-energy tensor (or \(T^{00}/c^2\)
  in the static limit), serving as the covariant scalar source.
\item
  \(a_0 = c\,H_0/(2\pi)\) from Theorem T20.
\item
  \(G = c^3\ell_P^2/\hbar\) from Theorem T19.
\end{itemize}

\textbf{Structural status.} Form FORCED at substrate level by
Arc\_ED10\_4 candidate-form selection chain: scalar-tensor
acoustic-metric uniquely substrate-compatible; \(f(R)\) inadmissible;
bimetric collapses to scalar-tensor. Weak-field limit reduces to Section
7's modified Poisson equation under three load-bearing structural
assumptions. OS positivity preserved + ghost-free + gradient-stable
under structural conditions; deep-MOND superluminality structurally
inherited.

\textbf{This is not GR.} No Einstein-Hilbert kinetic term for the
metric; no dynamical-metric-as-fundamental-field content. The metric is
acoustic-class kinematic; the equation is a substrate-derived
covariantization of the MOND-class modified Poisson equation,
structurally equivalent to the standard RAQUAL / scalar-tensor MOND
covariantizations in the literature with the ED-program-specific
contribution being the substrate-derivation chain.

\begin{center}\rule{0.5\linewidth}{0.5pt}\end{center}

\subsection{10. Structural Verdict}\label{structural-verdict}

\subsubsection{What is FORCED at substrate
level}\label{what-is-forced-at-substrate-level}

\begin{itemize}
\tightlist
\item
  Newton's gravitational coupling form \(G = c^3\ell_P^2/\hbar\) (T19).
\item
  Transition acceleration form \(a_0 = c\,H_0/(2\pi)\) +
  \(\ell_P\)-invariance under continuum limit (T20).
\item
  ED Combination Rule \(a = \sqrt{a_N\,a_0}\) (ECR).
\item
  BTFR slope-4 relation \(v^4 = G\,M\,a_0\) (T21).
\item
  Modified Poisson equation form + asymptotic constraints on \(\mu(x)\)
  (Section 7).
\item
  Substrate cumulative-strain four-index → Riemann-class object (Section
  8).
\item
  Acoustic-metric + Christoffel-class connection + Riemann-like
  curvature mappings (Section 8).
\item
  Scalar-tensor acoustic-metric covariantization as the unique
  substrate-compatible form (Section 9).
\item
  Weak-field limit reduction of covariant equation to flat-background
  modified Poisson (Section 8).
\item
  OS positivity + ghost-freedom + gradient stability under structural
  conditions (Section 8).
\end{itemize}

\subsubsection{What is INHERITED at value
layer}\label{what-is-inherited-at-value-layer}

\begin{itemize}
\tightlist
\item
  Numerical values of \(G\) and \(a_0\) (via \(\ell_P\) and \(H_0\)).
\item
  Detailed shape of \(\mu(x)\) in transition regime.
\item
  Curvature-coupling coefficients \(\alpha_R\), \(\beta_{R_{\mu\nu}}\)
  at value layer.
\item
  Cosmological boundary conditions; specific mass distributions.
\item
  Specific compact gauge group choice (for the matter sector via T17
  minimal coupling).
\end{itemize}

\subsubsection{What is structurally
conditional}\label{what-is-structurally-conditional}

\begin{itemize}
\tightlist
\item
  \textbf{Deep-MOND superluminality} is structurally FORCED in the
  deep-MOND regime as the structural cost of the substrate-derivation
  chain. Avoidance routes structurally inadmissible.
\item
  \textbf{Curvature-coupling coefficients must remain subleading} for
  the weak-field limit to reduce cleanly to SG-4. Per closed-arc work;
  not structurally proven within Arc ED-10.
\item
  \textbf{Acoustic metric must remain Lorentzian and non-degenerate.}
  FORCED in non-extreme regimes per ED-Phys-10; could fail near extreme
  substrate-density configurations (out of scope).
\end{itemize}

\subsubsection{What is not obtained}\label{what-is-not-obtained}

\begin{itemize}
\tightlist
\item
  \textbf{No general relativity emergence.} The Einstein equation
  \(G_{\mu\nu} = 8\pi G T_{\mu\nu}\) is not derived; remains SPECULATIVE
  per ED-Phys-10 acoustic-metric-only baseline.
\item
  \textbf{No metric dynamics beyond kinematic content.} Gravitational
  waves, black-hole solutions, cosmological metric evolution out of
  scope.
\item
  \textbf{No prediction of numerical \(G\), \(a_0\), or \(H_0\) values.}
  Numerical content INHERITED throughout.
\item
  \textbf{No solution of any Clay Millennium Problem.} Curvature
  emergence is not itself a Clay problem.
\item
  \textbf{No constructive-rigorous OS reconstruction theorem.}
  OS-positivity audit is structural-suggestive level only.
\end{itemize}

The verdict is \textbf{conditional-positive at structural-suggestive
level} parallel in form and honesty to the program's NS-Smoothness
Intermediate Path C verdict for Clay-NS and the Yang-Mills
Clay-relevance verdict.

\begin{center}\rule{0.5\linewidth}{0.5pt}\end{center}

\subsection{11. Discussion and Outlook}\label{discussion-and-outlook}

\subsubsection{Relation to MOND
literature}\label{relation-to-mond-literature}

The substrate-derived modified Poisson equation (Section 7) is
structurally equivalent to the standard MOND-class AQUAL field equation
introduced by Bekenstein-Milgrom 1984. The substrate-derived
scalar-tensor covariantization (Section 9) is structurally equivalent to
RAQUAL / scalar-tensor MOND covariantizations studied since the same
period. The ED-program-specific contribution is \textbf{not} novel
field-equation content; it is the \textbf{substrate-derivation chain}
identifying these equations as substrate-FORCED rather than
phenomenologically chosen. The originality lies in the
substrate-grounding of form-selection, not in the form itself.

This positions ED's substrate-gravity work as a \emph{foundational}
contribution to the MOND-class research program rather than as an
additional phenomenological MOND-variant. Whether MOND-class behavior is
the correct phenomenological framework for galactic dynamics +
cosmological-scale gravity is an empirical question that the ED
substrate-derivation does not by itself resolve; what the ED
substrate-derivation supplies is a structural reason \emph{why}
MOND-class equations could be expected to emerge from a deeper
structural framework.

\subsubsection{Relation to analog
gravity}\label{relation-to-analog-gravity}

The acoustic-metric reading of curvature emergence (Section 8) connects
to the analog-gravity research program studying acoustic / hydrodynamic
/ condensed-matter analogs of gravitational physics. In the standard
analog-gravity framework, the metric emerges as the effective metric
experienced by acoustic perturbations propagating through a flowing
fluid; the metric is kinematic in the sense identified here. ED's
substrate-derivation supplies an analog-gravity-class metric-emergence
at the deeper substrate level, with the difference that ED's substrate
is not a fluid in the standard hydrodynamic sense but a
participation-channel network. The acoustic-metric emergence is
structurally common; the substrate identification is ED-specific.

\subsubsection{Future arcs}\label{future-arcs}

Several future arcs could extend the substrate-gravity content:

\begin{itemize}
\tightlist
\item
  \textbf{Substrate Gravity Phase-2.} Optional extensions within the
  flat-background regime --- dynamic mass-distribution-evolution
  effects, cosmological-perturbation-class extensions, explicit MOND
  empirical-fit work. Discretionary scope.
\item
  \textbf{Cosmological extensions.} The transition acceleration
  \(a_0 \propto H_0\) supplies a cosmological-scale connection that
  could be developed further at the level of substrate-gravity-cosmology
  coupling. Out of scope of the present paper but a natural extension
  direction.
\item
  \textbf{Stronger curvature regimes.} If ED-Phys-10's
  acoustic-metric-only baseline is revisited, the curvature-emergent
  framework could potentially extend to genuinely curvature-dominated
  regimes (black-hole-class substrate configurations,
  gravitational-wave-class metric perturbations). This would require
  structural commitment beyond the present arc and is not undertaken
  here.
\end{itemize}

\subsubsection{ED Program Manifesto}\label{ed-program-manifesto}

The substrate-gravity content presented here is one sector of the
broader Event Density program. The full program covers Phase-1 QM
emergence, the closed structural-foundation theorems (T17 / T18 / N1 /
GR1 / DCGT), the NS / MHD applied-arc program with NS Synthesis Paper,
the Yang-Mills arc, and the present substrate-gravity work. A
long-horizon ED Program Manifesto would consolidate all of this into a
single book-length program-overview publication. The present
substrate-gravity-foundations paper + the existing NS Synthesis Paper +
the ED-QFT Unified Overview Paper form the publication-grade backbone
for that future Manifesto.

\begin{center}\rule{0.5\linewidth}{0.5pt}\end{center}

\subsection{References}\label{references}

(Placeholder for later insertion. Anticipated references include:)

\begin{itemize}
\tightlist
\item
  \emph{Event Density: One Substrate, Three Domains --- A Structural
  Account of Quantum Mechanics, Galactic Gravity, and Universal
  Soft-Matter Mobility} (Proxmire 2026). Earlier program-overview.
\item
  \emph{The Architectural Foundations of Navier-Stokes in Event Density}
  (Proxmire 2026). NS / MHD / YM closed-arc reference paper.
\item
  \emph{The Event-Density Foundations of Quantum Field Theory: A Unified
  Substrate-to-Continuum Overview} (Proxmire 2026). Program-level
  synthesis.
\item
  \emph{The Primitive-Level Arrow of Time in Event Density: The Closure
  of Arc B and Theorem 18} (Proxmire 2026). T18 paper.
\item
  \emph{Gauge Fields as Forced Rule-Type Structure: Theorem 17}
  (Proxmire 2026). Arc Q closure paper.
\item
  \emph{Fields and Forces in Event Density} (Proxmire, Feb 2026,
  ED-I-06). Ontological framework paper.
\item
  Bekenstein, J.D. and Milgrom, M. (1984). Does the missing mass problem
  signal the breakdown of Newtonian gravity? Astrophys. J. 286.
\item
  Milgrom, M. (1983). A modification of the Newtonian dynamics as a
  possible alternative to the hidden mass hypothesis. Astrophys. J. 270.
\item
  Bekenstein, J.D. (2004). Relativistic gravitation theory for the
  modified Newtonian dynamics paradigm. Phys. Rev.~D 70.
\item
  Sanders, R.H. and McGaugh, S.S. (2002). Modified Newtonian dynamics as
  an alternative to dark matter. Annu. Rev.~Astron. Astrophys. 40.
\item
  Lelli, F., McGaugh, S.S., Schombert, J.M. (2016). SPARC: Mass models
  for 175 disk galaxies with Spitzer photometry and accurate rotation
  curves. Astron. J. 152.
\end{itemize}

\begin{center}\rule{0.5\linewidth}{0.5pt}\end{center}

\emph{Foundations of Substrate Gravity: From Newtonian Limit to
Curvature Emergence in the Event Density Framework. Canonical
substrate-gravity-and-curvature-emergence reference for the ED program.
Integrates T19 (Newton's \(G\) from substrate) + T20 (\(a_0\)
horizon-scale + \(\ell_P\)-invariant) + ECR (geometric-mean composition)
+ T21 (BTFR slope-4) + Arc SG (flat-background modified Poisson
equation) + Arc ED-10 (scalar-tensor acoustic-metric covariantization)
into a single publication-grade synthesis. Form-FORCED / value-INHERITED
methodology consistent across every section. Conditional-positive
structural verdict parallel in form and honesty to NS-Smoothness
Intermediate Path C and Yang-Mills Clay-relevance verdicts. No claim of
GR derivation; ED-Phys-10 acoustic-metric guardrails preserved.
Substrate-derivation chain identifies MOND-class field equation +
scalar-tensor covariantization as substrate-FORCED + supplies structural
reason for BTFR slope-4 + identifies deep-MOND superluminality as
structural cost.}

\end{document}
